\documentclass[10pt,a4paper]{article}
\usepackage{amsmath,amssymb,amsthm,graphicx}
\usepackage{hyperref}
\usepackage{latexsym}
\usepackage{amssymb}
\usepackage{amsthm}
\usepackage{ulem}
\usepackage{enumitem}
\setlist{itemsep=0.4pt,topsep=0pt}
\parindent 0pt
\parindent 0pt
\parskip 12pt
\newcommand{\HCF}{\mathrm{hcf}}
\newcommand{\ZZ}{\mathbb{Z}}
\newcommand{\rk}{\mathrm{rk}}
\DeclareMathOperator{\len}{length}
\DeclareMathOperator{\bl}{b}
\DeclareMathOperator{\cl}{cl}
\DeclareMathOperator{\di}{d}
\DeclareMathOperator{\Cl}{Cl}
\DeclareMathOperator{\acl}{acl}
\DeclareMathOperator{\tp}{tp}
\DeclareMathOperator{\Th}{Th}
\newcommand{\Cor}{\mathcal{C}_\mathrm{or}}
\newcommand{\Cun}{\mathcal{C}_\mathrm{un}}
\newcommand{\mons}{\mathbb{M}}
\newcommand{\strong}{\mathrm{strong}}
\newcommand{\axioms}{\Delta_{\mathrm{FSTS}}}
\newtheoremstyle{note}% hnamei
{8pt}% hSpace abovei
{14pt}% hSpace belowi
{\normalfont}% hBody fonti
{}% hIndent amounti1
{\bfseries}% hTheorem head fonti
{ }% hPunctuation after theorem headi
{.5em}% hSpace after theorem headi2
{}% hTheorem head spec (can be left empty, meaning ‘normal’)i
\theoremstyle{note}
\newtheorem{thm}{Theorem}[section]
%\newtheorem{theorem}{Theorem}
\newtheorem{lemma}[thm]{Lemma} 
\newtheorem{df}[thm]{Definition}
\newtheorem{rmk}[thm]{Remark} 
\newtheorem{prop}[thm]{Proposition}
\newtheorem{fact}[thm]{Fact}
\newtheorem{cor}[thm]{Corollary}
\newtheorem{conj}[thm]{Conjecture}
\newtheorem{ex}[thm]{Example}
\newtheorem{qu}[thm]{Question}
\newtheorem*{notation}{Notation} 
\title{The class of unconfined systems has the bounded amalgamation property with an explicit bound}
\author{Silvia Barbina, Enrique Casanovas}
\date{arxiv:2411.13723v3; Creative Commons Attribution 4.0 International}
\begin{document}
\maketitle
\noindent\textit{Editorial note.} The counted statement is the author's environment at source0.tex lines 740-741 and its proof environment opens at line 742; the pairing is by adjacency and the retained proof range is the whole environment.\par\medskip
\noindent\textit{Editorial note.} The preamble is the author's own package lines, macro block and theorem-environment declarations (source lines 4-65); the author's optioned amsmath line and the fleqn-dependent layout line are not retained.\par\medskip
\noindent\textit{Editorial note.} The source carries no label-number dump and no counter values are imposed, so the printed numbering is the wrapper's default; identity is the author label, or the pinned source position where the author gave none.\par\medskip
\noindent\textit{Editorial note.} Metadata provenance: the source record of the pinned copy and the licence evidence file are both present; the licence is CC BY 4.0 with status supported, and the card's source\_version matches the value recorded in that evidence file.\par\medskip
\noindent\textit{Editorial note.} r002 note: r001 failed because the retained author proof cites two bibliography keys of the article; r002 bundles the article's own bibliography file and changes nothing else.\par\medskip
\noindent\textit{Editorial scope.} One original result of the pinned author manuscript, counted once: the bounded amalgamation property with an explicit bound for unconfined systems (author label Feb\_23\_2024\_2).\par\medskip
\noindent\textit{Allowed standard prerequisites (editorial).}\begin{itemize}
\item Amalgamation in a bound. A\$\textbackslash{}leq\$n B means that the extension absorbs any set of at most n new elements.. It is the bound that the proposition makes explicit.
\item Two HF-orderings combined. An HF-ordering of C can be concatenated with the restriction of an HF-ordering of B* to the complement of C to give an HF-ordering of the union.. It is the construction that produces the amalgam.
\item Isomorphic copy over A. B* is a copy of B fixing A, arranged so that the union with C is unconfined.. It is the choice made at the start of the proof.
\end{itemize}\medskip
\begin{lemma}\label{Feb_23_2024_1}\begin{enumerate}
			\item Assume $A\leq B$, $ b\in B\smallsetminus A$ and there is a block $\{a_1,a_2,b\}$  with  $a_1,a_2\in A$. Then  $Ab\leq B$.
			\item Assume $A\leq^n B$, $b\in B\smallsetminus A$ and there is a block $\{a_1,a_2,b\}$  with  $a_1,a_2\in A$. Then  $Ab\leq^{n-1} B$.
		\end{enumerate}
	\end{lemma}

	\begin{prop}\label{Feb_23_2024_2} $(\Cun,\leq)$ has the $n$-amalgamation property for every $n$.  In fact, if $A\leq B$ and $n<\omega$ are given, we can always take $m= n+|B\smallsetminus A|$.
	\end{prop}

	\begin{proof} Assume $A\leq B$  and  $A\leq^m C$, where $m= n+|B\smallsetminus A|$. We want to find $B^\ast\cong_A B$ and $D\in \Cun$  such that  $C\leq D$ and $B^\ast\leq^n D$.  We may assume $C\cap B=A$. Note that the blocks containing only elements of $A$ are the same in $B$ and in $C$.  Assume there are $b\in B\smallsetminus A$, $c\in C\smallsetminus A$ and  $a_1,a_2\in A$  such that $\{a_1,a_2,b\}$ is a block of $B$ and  $\{a_1,a_2,c\}$ is a block of $C$. Redefine $A^\prime= A\cup\{c\}$,  $B^\prime = (B\smallsetminus \{b\})\cup \{c\}$ and modify the blocks of $B$ by exchanging $b$ with $c$  to obtain the blocks of $B^\prime$. Then $B^\prime\cong_A B$ and $B^\prime\cap C=A^\prime$.  By Lemma~\ref{Feb_23_2024_1}, we have  $A^\prime\leq B^\prime$ and $A^\prime\leq^{m-1} C$. Iterating this process gives $B^\ast\cong_A B$  such that if $A^\ast =B^\ast \cap C$, then   $A^\ast \leq B^\ast$,  $A^\ast\leq^{n +i}C$  for some $i \geq 0$  and there are no $b\in B^\ast\smallsetminus A^\ast$, $c\in C\smallsetminus A^\ast$ and $a_1,a_2\in A^\ast$ such that  $\{a_1,a_2,b\}$ is a block of $B^\ast$  and  $\{a_1,a_2,c\}$ is a block of $C$.  Now consider $D=B^\ast\cup C$ as a free amalgam of $B^\ast$ and $C$ over $A^\ast$, that is,  with the blocks in $B^\ast$ and in $C$ and no others.  Then $D$ is a partial STS. We claim that $D\in \Cun$, $C\leq D$ and $B^\ast\leq^n D$.
		
First we show that $D$ has an HF-ordering where $C$ is an initial segment. This proves that $D\in\Cun$ and $C\leq D$.  Take an HF-ordering $<_C$ of $C$ and an HF-ordering $<_{B^\ast}$ of  $B^\ast=A^\ast\cup (B^\ast\smallsetminus C)$ where $A^\ast$ is an initial segment. Then the sum 
$$(C,<_C)+(B^\ast\smallsetminus C,<_{B^\ast}\restriction (B^\ast\smallsetminus C))$$ 
is an HF-ordering of $D$. If not, there is an element of $B^\ast\smallsetminus C$ that belongs to different blocks containing smaller elements, but these elements then are in $A^\ast\cup (B^\ast\smallsetminus C)$, contradicting the choice of the ordering $<_{B^\ast}$ of $B^\ast$ as an HF-ordering over $A^\ast$.
		
		Finally, we check that $B^\ast\leq^n D$.  Let $X\subseteq D\smallsetminus B^\ast = C\smallsetminus A^\ast$  with $|X|\leq n$; we show that $B^\ast\leq B^\ast \cup X$.  Since  $A^\ast\leq B^\ast$, we have an HF-ordering $<_{B^\ast}$ of $B^\ast$ with initial segment $A^\ast$.  Since $A^\ast\leq^n C$, we know that $A^\ast\leq A^\ast\cup X$, and hence there is an HF-ordering $<_{A^\ast \cup X}$ of $A^\ast\cup  X$ with initial segment $A^\ast$. We combine these orderings by taking the sum $(B^\ast,<_{B^\ast})+(X,<_{A^\ast\cup X}\restriction X)$. Since an element of $X$ cannot have a block containing elements of $B^\ast$ unless the elements belong to $A^\ast$, we obtain an HF-ordering of $B^\ast\cup X$ with initial segment $B^\ast$.
	\end{proof}

\begin{thebibliography}{9}
\bibitem{barcas2}
S.~Barbina and E.~Casanovas.
\newblock Endomorphisms of free {S}teiner quasigroups, 2024.
\newblock{\url{arXiv:2412.06055}}.

\bibitem{casanovas}
E.~Casanovas.
\newblock Compactly expandable models and stability.
\newblock {\em The Journal of Symbolic Logic}, 60:673--683, 1995.
\bibitem{ganter1}
B.~Ganter and H.~Werner.
\newblock Equational classes of {S}teiner systems.
\newblock {\em Algebra Universalis}, 5:125--140, 1975.

\bibitem{ganter2}
B.~Ganter and H.~Werner.
\newblock Co-ordinatizing {S}teiner systems.
\newblock In C.C. Linder and A.~Rosa, editors, {\em Topics on {S}teiner
  systems}, pages 3--24. North Holland, 1980.
\end{thebibliography}
\end{document}
